%!TEX root=./neurips_2023.tex

\section{Background}
\label{sec:background}

% LoRA & IA3
% Before introducing merge approaches, we should first clarify what are we merging when referring to parameter-efficient modules.
    % For any model layer expressed as a matrix multiplication of the form h=Wx, adding parameter efficient modules results in a additional $\Delta h$ to it.
    % LoRA injects trainable low-rank decomposition matrices A & B into the layers of a pretrained model. 
    % (IA)^3 introduces trainable vectors l into different components of a Transformer model which perform element-wise rescaling of inner model activations. 
    % when we are performing arithmeric operations between parameters of these modules, what we are changing is the Delta h in the whole function. Therefore what we are doing is not simply add them together or doing substraction between them, we focused on the value of h and delta h.

Parameter-efficient finetuning was first introduced by~\citet{houlsby2019parameter} into NLP, where they propose to insert small modules called adapters into the pretrained transformer~\citep{attention_transformer} at different places, such as after the attention module and after the feed-forward module within each layer. 
During finetuning, only the adapter parameters are updated.
% The adapter layer takes an input vector $\vh \in \sR^d$ and modifies it as:
% \begin{equation}
% \label{eq:adapter}
%     \vh \leftarrow \vh + f(\vh\wdd)\wu,
% \end{equation} 
% where $\wdd \in \sR^{d\times r}$ is the down-projection matrix and $\wu\in \sR^{r\times d}$ is the up-projection matrix. 
% Parameters of these two projection matrices are all the learnable parameters in adapters.
% $r \ll d $ is the bottleneck dimension, and $f$ is a nonlinear activation function like ReLU. 
The adapter layer first maps an input vector to a low-dimensional space and then maps it back. 
This bottleneck projection architecture is widely adopted in later work~\citep{pfeiffer2021adapterfusion,NEURIPS2021_081be9fd,hu2021lora},
and~\citet{he2021towards} show that many PEFT methods could be viewed as a form of adapter. 
In this paper, we focus on two recent state-of-the-art PEFT methods, LoRA~\citep{hu2021lora} and \ia~\citep{liu2022few}, which we describe below.

\paragraph{LoRA}
is probably the most effective PEFT method to date given its superior performance as reported in~\citet{hu2021lora}. 
It has notably garnered increasing interest recently, becoming a standard approach for adapting large language models such as LLaMA~\citep{touvron2023llama} under limited computational resources~\citep{alpacalora}. 
LoRA bears a similar form to adapter, albeit with minor differences. Specifically, for any weight matrices in the transformer that take an input $\vx \in \sR^k$ and output $\vh \in \sR^d$, LoRA modifies $\vh$ as:
\begin{equation}
    \label{eq:lora}
    \vh \leftarrow \vh + \mB\mA\vx,
\end{equation}
where $\mB \in \sR^{d\times r}, \mA \in \sR^{r\times k}$ are the projection matrices, and the rank $r \ll \min(d, k)$.
While LoRA could be applied for any weight matrices,~\citet{hu2021lora} utilize it in the query and value projection matrices of the attention module practically. In this study, we adhere to this established practice. 
In LoRA tuning, $\mA$ is initialized following random Gaussian distribution, and $\mB$ is initialized to all zeros to recover the pretrained model at the beginning, as suggested by~\citet{hu2021lora}. 
$\tl=\{\mA, \mB\}$ forms the parameter-efficient module in LoRA, which we aim to compose with other LoRA modules trained differently.
% For any model layer expressed as $h=Wx$ where $W\in \mathbb{R}^{d\times k}$, the integration of trained PEMs results in an additional $\Delta h$. For LoRA~\citep{hu2021lora}, trainable low-rank decomposition matrices ($A$ and $B$) are injected into the layers of a pretrained model. The hyperparameter $r$ determines the low-dimensional rank of the decomposition, with $A$ having dimensions $r\times k$ and $B$ having dimensions $d\times r$.

% \begin{equation}
%     h_{\rm{LoRA}}=Wx+\frac{\alpha}{r}BAx.
% \end{equation}

% Therefore, the resulting $\Delta h$ for LoRA is:

% \begin{equation}
% \Delta h_{\rm{LoRA}}=\frac{\alpha}{r}BAx
% \end{equation}

% where A is initialized by a random Gaussian initialization and B by a zero matrix.
\paragraph{\ia} is proposed by~\citet{liu2022few} for few-shot learning. 
It introduces trainable vectors $\vl_k$, $\vl_v$, and $\vl_{ff}$ to respectively rescale the attention keys, attention values and the inner activations in position-wise feed-forward networks. 
Let the modified hidden states be $\vh$, \ia modifies it as:
\begin{equation}
\label{eq:ia3}
    \vh \leftarrow \vl \odot \vh,
\end{equation}
where $\vl$ are initialized as all ones so that the model is unchanged at the beginning of tuning. $\ti=\{\vl_k, \vl_v, \vl_{ff}\}$ form the PEM of \ia~that we aim to compose.

% In the case of \ia~\citep{liu2022few}, trainable vectors ($l$) are introduced to rescale the inner model activations element-wise within a Transformer model.

% \begin{equation}
% h_{\rm{(IA)^3}}=l \odot Wx.
% \end{equation}

% Thus, the resulting $\Delta h$ for \ia~ is:

% \begin{equation}
% \Delta h_{\rm{(IA)^3}}=(l-1)\odot Wx
% \end{equation}

% where $l$ in \ia~ is initialized using an all-ones vector. 



\section{Composition through Arithmetic Operation}
\label{sec:definition}
Prior work compose PEMs trained on different tasks for multi-task purposes through learning to fuse their outputs~\citep{pfeiffer2021adapterfusion,wang2022adamix}.
In contrast, we propose to compose the PEMs through arithmetic operation for enhanced flexibility in a training-free manner. 
Our method is inspired by recent study on the linear connectivity of trained models in a full finetuning setting~\citep{wortsman2022model,matena2022merging,ainsworth2022git,jin2022dataless}.
These studies suggest that parameters of tuned models can be directly added to improve generalization, provided they are initialized from the same pretrained model checkpoint.
The underlying hypothesis is that two models finetuned from the same pretrained checkpoint often lie in the same error basin~\citep{NEURIPS2020_0607f4c7}, and thus the parameters could be directly added. 
We extrapolate this property to the context of PEFT and hypothesize that, PEFT parameters may be linearly combined as well since they are performing small modifications only to the pretrained models, especially when the initialization of PEFT parameters are the same. 
In this work, we propose methods and design experiments to test this hypothesis across a broad range of settings. 
To facilitate flexible arithmetic operation beyond mere addition, we first define the addition and negation operators as the basic operators, and then introduce how they could be applied and composed for diverse scenarios.

% Performing linear arithmetic operations on the parameters of PEMs in the weight space alters the resulting $\Delta h$ in various ways, therefore specialized definition is necessary.
% We define the addition operator as $\oplus$ and negation as $\ominus$ for PEMs, which act differently depending on the PEFT methods.
% taking into account both direct arithmetic operations on the parameters and the resulting values of $h$ and $\Delta h$.

\subsection{Basic Operators}
% We define two basic operators, the addition and the negation operator for both $\tl$ and $\ti$, as detailed below.

% \subsection{Composing Operators}
% Consider the two trained PEMs with parameters represented as $(A_1,B_1)$ and $(A_2,B_2)$ in LoRA~\citep{hu2021lora} as well as $l_1$ and $l_2$ in \ia~\citep{liu2022few}. 
% To ensure that both modules are inside the same loss minima basin as much as possible after training, identical pretrained models, module architectures, and module initialization are used throughout our experiments to enhance mode connectivity. 
% Additionally, we introduce $\lambda \in (0,1)$ to allocate proper proportions to merged modules.


% We vary lambda from 0 to 1 with a step size of 0.02 to optimize the merge setting.(move to exps section)

\paragraph{PEM addition operator:} 
Similar to previous work on linearly combining parameters, 
we define module addition as the operation of pairing the arguments at corresponding positions and adding them component-wise. 
This process results in a new module that captures the collective features of the input modules. 
Formally, for parameters of two PEMs $\ta$ and $\tb$, we define the addition operator $\oplus$ as:
\begin{equation}
\label{eq:add}
    \tp = \ta \oplus \tb = \ta + \tb,
\end{equation}
where we use $\vtheta$ to represent PEM parameters in general, and $\tp$ represents the merged parameters. 
Eq.~\ref{eq:add} applies to both $\tl$ and $\ti$.
% For LoRA, the addition of two trained modules in proportion $\lambda$ and $(1-\lambda)$ respectively results in:

% \begin{equation}
% B_{add}=\lambda B_1+(1-\lambda) B_2,\ \text{and} \ A_{add}=\lambda A_1+(1-\lambda) A_2.
% \end{equation}

% Similarly, for $\rm{(IA)^3}$, the resulting module addition is:

% \begin{equation}
% l_{add}=\lambda l_1+(1-\lambda)l_2.
% \end{equation}

% Put them into Appendix
% \begin{equation}
% \Delta h_{\rm{LoRA}\_add}=\frac{\alpha}{r}[\lambda B_1+(1-\lambda) B_2][\lambda A_1+(1-\lambda) A_2]x.
% \end{equation}

% \begin{equation}
% \Delta h_{\rm{(IA)^3}\_add}=[\lambda l_1+(1-\lambda)l_2-1]\odot Wx.
% \end{equation}

\paragraph{PEM negation operator:}
The objective of the negation operator is to facilitate unlearning or forgetting certain skills, for example, a PEM trained on toxic data may be directly negated as a plug-in detoxifier. 
With the predefined addition operator, the negation operator $\ominus$ could naturally enable the subtraction operation as $\ta\ominus\tb = \ta\oplus(\ominus\tb)$. 
% to observe a decrease in the skill level of the model via negation, resulting in obtaining the opposite value of $\Delta h$. 
% While the addition operator is pretty straightforward to define, the negation operator is more complicated and needs to be considered case by case with respect to different PEMs -- for example, one cannot simply take negation of all parameters of LoRA, which would essentially change nothing. 
Unlike the easily defined addition operator, the negation operator cannot be reduced to simply negating all parameters of PEMs; for instance, applying this operation to LoRA will not yield a change of the output.
To properly formulate the negation operator, we focus on the modification that the PEMs apply to the hidden states $\vh$.
The intuition is that we can view all PEFT methods as applying a modification $\Delta \vh$ added to the original $\vh$, which is a general and unified perspective to view PEFT methods as proposed in~\citet{he2021towards}. 
Since $\Delta \vh$ is adding certain skills to the model hidden states, and we propose to design PEM negation operator to negate $\Delta \vh$: 
% \jz{I want to express negation as $-\Delta \vh$ and show its instantiation for LoRA and \ia here.}\jh{please put a word ``negate'' on top of the arrow}
\begin{equation}
    \vh \leftarrow \vh + \Delta \vh \ 
    \xRightarrow{\text{negate}} \ \vh \leftarrow \vh + (-\Delta \vh)
\end{equation}
Specifically, for LoRA and \ia~we have:
\begin{equation}
    \label{eq:delta}
    \Delta \vh_{\text{lora}} = \mB\mA\vx, \quad \Delta \vh_{\text{ia3}} = (\vl - \mathbf{1}) \odot \vh_{\text{ia3}},
\end{equation}  
then to negate $\Delta \vh_{\text{lora}}$, we could simply negate $\mB$ or $\mA$ while keeping the other unchanged. 
Practically in our experiment, we choose to negate $\mB$ as:
% \footnote{In Appendix~\ref{sec:apx_multicombination}\jh{check}\jz{analogy} we conduct analysis on negating $\mA$ instead, observing that the results are robust to the choice of the negation matrix.}
\begin{equation}
\tnl = \ominus\tl = \{\mA, -\mB\}.
\end{equation}
For a specified $\vl$ vector in \ia, we solve the equation on negating $\Delta \vh_{\text{ia3}}$ and obtain:
\begin{equation}
\label{eq:ia3-neg}
    (\vl^{\text{neg}} - \mathbf{1}) \odot \vh_{\text{ia3}} = -(\vl - \mathbf{1}) \odot \vh_{\text{ia3}} \ \Rightarrow \ \vl^{\text{neg}} = \ominus \vl = \mathbf{2}-\vl.
\end{equation}
Eq.~\ref{eq:ia3-neg} is applied to all the three $\vl$ vectors to negate the \ia~module. 
We also include an ablation analysis on negation operator for both LoRA and \ia~ in Appendix \ref{sec:apx_example_toxic}.
Next, we demonstrate how to utilize the two basic operators $\oplus$ and $\ominus$ in different scenarios.

\begin{table}[]
\centering
\caption{Different settings studied in this work and their corresponding arithmetic operations.}
\label{table:expls}
\begin{tabular}{lll}
\toprule
Settings & Arithmetic operations \\ \midrule
% \ref{sec:datasetsplit} & 
Distribution generalization & $\ta \oplus \tb$ \\
% \ref{sec:multitask} & 
Multi-tasking & $\ta \oplus \tb$ \\
% \ref{sec:neg} & 
Unlearning & $\ominus \vtheta$ \\
% \ref{sec:combine} & 
Domain transfer & $\ta \ominus \tb \oplus \tc$ \\
Detoxifying instruction-tuned LLMs & $\ta \ominus \tb$ \\ 
% \ref{sec:llama} & 
\bottomrule
\end{tabular}
\end{table}

\subsection{Composing Basic Operators}
When we apply the basic operators to merge different PEMs in practice, a weight hyperparameter $\lambda\in[0,1]$ is required to alter the relative weights of the modules, as in~\citet{ilharco2022editing,wang2022adamix}. Therefore, we compute $\tp$ as a linear interpolation of two modules and assign a weight scalar to $\tn$ as follows:
\begin{equation}
    \tp = \lambda \ta \oplus (1-\lambda) \tb, \quad \tn=\ominus \lambda\vtheta.
\end{equation}
$\lambda$ is a hyperparameter that is tuned on a validation set. 
While advanced methods of reweighting different parameters in the full finetuning setting have been proposed by~\citet{matena2022merging,jin2022dataless}, we leave exploration on this aspect as future work and focus on the simplest version in this paper.
Our empirical study next covers four different arithmetic operations based on the operators, as listed in Table~\ref{table:expls}:\footnote{Here we omit the $\lambda$ hyperparameter for ease of notations.} (1) $\ta \oplus \tb$ for distribution generalization or multi-task learning; (2) $\ominus \vtheta$ for unlearning certain abilities from a pretrained model; (3) $\ta \ominus \tb \oplus \tc$ for transferring a model across domains -- for example, when $\ta$ represents a classification model trained on restaurant reviews, $\tb$ denotes a language model on restaurant reviews, and $\tc$ signifies a language model on product reviews, then $\ta \ominus \tb \oplus \tc$ may lead to a PEM for classification on product reviews. Such an analogy computation resembles the well-known word embedding example ``\texttt{queen = king - man + woman}'', and has been verified in a full finetuning setting by~\citet{ilharco2022editing}; and (4) $\ta \ominus \tb$ for detoxifying instruction-tuned LLMs. 
% \jh{maybe a table}


% We highlight that negation is a complicated process that cannot be achieved simply by taking the negative of all parameters within the module.
% For LoRA $\theta$, the negated module $\theta_{neg}$ is defined as $\ominus \theta$, where $B_{neg}=-B$ and $A_{neg}=A$.

% \begin{equation}
% B_{neg}=-B,\ \text{and} \ A_{neg}=A.
% \end{equation}

% For \ia~, to obtain the negated module, we calculate the resulting value of $h_{\rm{(IA)^3} \_ neg}$ according to:

% \begin{equation}
% h_{\rm{(IA)^3}\_neg}=h_{\rm{(IA)^3}}-\Delta h_{\rm{(IA)^3}}=(2-l)\odot Wx=l_{neg}\odot Wx.
% \end{equation}

% Therefore, $l_{neg}$ is given by:

% \begin{equation}
% l_{neg}=2-l.
% \end{equation}

% % multiply every parameter in lora_B by -1
% \begin{equation}
% \Delta h_{\rm{LoRA}\_neg}=\frac{\alpha}{r}(-B)Ax
% \end{equation}

% \textbf{Task analogy.}

% We take task analogy as an example to show the combination of PEMs through linear arithmetic operation incorporating the aforementioned definition.
% Consider the four modules $\theta_A$, $\theta_B$, $\theta_C$ and $\theta_D$, which have been trained to carry out tasks of the form "A is to B as C is to D". 
% By assembling modules $\theta_A$, $\theta_B$ and $\theta_C$, we derive the parameters of module $\theta_D$ expressed as:

% \begin{equation}
% \theta_D=\lambda \theta_B \oplus (1-\lambda)[\theta_C \ominus \theta_A)]
% \end{equation}

% \jh{be more concrete, what would be shown in appendix A?}Further information can be found in the Appendix \ref{sec:apx_formula}.

% % substraction
% \begin{equation}
% \Delta h_1-\Delta h_2=\Delta h_1 + (-\Delta h_2)
% \end{equation}

% \begin{equation}
% \Delta h_{\rm{LoRA}\_analogy}=\frac{\alpha}{r}[\lambda B_1-(1-\lambda) B_2][\lambda A_1+(1-\lambda) A_2]x
% \end{equation}

% % therefore analogy
% % lora
% \begin{equation}
% B_{\rm{LoRA}\_analogy}=\lambda B_1+(1-\lambda) (B_2-B_3)
% \end{equation}

% \begin{equation}
% A_{\rm{LoRA}\_analogy}=\lambda B_1+(1-\lambda) (A_2+A_3)
% \end{equation}

% % ia3
% \begin{equation}
% \Delta h_{\rm{(IA)^3}\_analogy}=[\lambda (l_1-1)+(1-\lambda)(l_2-1-l_3+1)]\odot Wx=[\lambda (l_1-1)+(1-\lambda)(l_2-l_3)]\odot Wx
% \end{equation}

% \begin{equation}
% h_{\rm{(IA)^3}\_analogy}=h_{\rm{(IA)^3}}-\Delta h_{\rm{(IA)^3}\_analogy}=[1-\lambda (l_1-1)-(1-\lambda)(l_2-l_3)]\odot Wx=l_{analogy}\odot Wx
% \end{equation}

% \begin{equation}
% l_{analogy}=1-\lambda (l_1-1)-(1-\lambda)(l_2-l_3)
% \end{equation}
